%%%%%%%%%%%%%%%%%%%%%%%%%%%%%%%%%%%%%%%%%%%%%%%%%%%%%%%%%%%%%%%%%%%%%%%%%%%
%%  chapter2.tex  –  Глава 2: Методы, алгоритмы и математические
%%                   основы для робастных динамических графовых нейронных сетей
%%  v9 – Единая постановка задачи с последовательным выводом.
%%       Декомпозиция на 7 подзадач, сопоставленных пробелам.
%%       Все доказательства, алгоритмы и анализ сложности сохранены.
%%       Только заголовки \section добавлены в оглавление.
%%%%%%%%%%%%%%%%%%%%%%%%%%%%%%%%%%%%%%%%%%%%%%%%%%%%%%%%%%%%%%%%%%%%%%%%%%%

\chapter{Методы, алгоритмы и математические основы для робастных динамических графовых нейронных сетей}
\label{ch:methods}

В настоящей главе разрабатывается полное теоретическое ядро диссертации. Глава начинается с единой постановки задачи, последовательно выводимой из фундаментальной задачи классификации графов путём поочерёдного включения каждого из четырёх требований (робастность, точность, эффективность, приватность) при трёх рабочих условиях (состязательная среда, нестационарность, распределённая динамика). Единая задача затем формально декомпозируется на семь подзадач, каждая из которых соответствует одному из пробелов, выявленных в Главе~\ref{ch:literature}. В оставшейся части главы представлены математическая формулировка, доказательства сходимости и робастности, алгоритмический псевдокод и анализ сложности для каждого из семи методов, решающих эти подзадачи. Глава завершается единой системой сертификации, связывающей все методы через регуляризацию согласованности.


%%=====================================================================
\section{Единая постановка задачи}
\label{sec:problem_statement}
%%=====================================================================

Последовательный вывод постановки задачи проходит через семь этапов, каждый из которых вводит новое математическое ограничение, соответствующее конкретному требованию или рабочему условию. Результатом является единая многоцелевая задача оптимизации, декомпозиция которой порождает семь методов, разработанных в настоящей диссертации.


\subsection{Этап 1: Классификация графов на динамических топологиях}

Рассмотрим темпоральный граф $\calG(t)=(\calV,\calE(t),\mathbf{X}(t))$, в котором узлы $\calV$, рёбра $\calE(t)$ и признаки $\mathbf{X}(t)=\{\bx_v(t)\in\R^{d}\}_{v\in\calV}$ эволюционируют во времени. Графовая нейронная сеть $f_\theta$ отображает граф в момент $T$ на предсказание. Базовая задача классификации:
\begin{equation}
\min_\theta\;
\E_{(\calG,y)\sim\calD}\!\bigl[
  \calL\!\bigl(f_\theta(\calG(T)),\,y\bigr)
\bigr],
\label{eq:ps_base}
\end{equation}
где $\calL$ --- кросс-энтропийные потери, а $\calD$ --- порождающее распределение данных. Данная формулировка предполагает, что граф чист, статичен при выводе, доступен централизованно, и что распределение данных не изменяется. Каждое из этих предположений нарушается на практике.


\subsection{Этап 2: Состязательная робастность}

В состязательной среде (Условие У1) целенаправленный противник возмущает граф. Включение Требования Т1 (робастность) преобразует задачу в формулировку мин-макс:
\begin{equation}
\min_\theta\;
\E_{(\calG,y)\sim\calD}\!\Bigl[
  \max_{\|\bdelta\|\le\epsilon_{\mathrm{adv}}}
    \calL\!\bigl(f_\theta(\calG(T)+\bdelta),\,y\bigr)
\Bigr],
\label{eq:ps_robust}
\end{equation}
где $\bdelta$ представляет возмущения признаков узлов, весов рёбер или топологической структуры, а $\epsilon_{\mathrm{adv}}$ --- бюджет возмущений. Это стандартная формулировка робастной оптимизации~\cite{Madry2018,Goodfellow2015,Szegedy2014}.


\subsection{Этап 3: Динамика непрерывного времени с сертифицированными границами}

Для получения сертифицированной робастности, а не эмпирического состязательного обучения, внутренняя динамика модели должна быть ограничена. Требование, чтобы $f_\theta$ эволюционировал через нейронное ОДУ с вектором поля, ограниченным по Липшицу, добавляет ограничение:
\begin{equation}
\frac{d\bh_v(t)}{dt} = f_\theta\!\bigl(\bh_v(t),\,\mathbf{A}(t),\,t\bigr),
\qquad
\|f_\theta(\bh_1,t)-f_\theta(\bh_2,t)\| \le L_g\|\bh_1-\bh_2\|,
\label{eq:ps_ode}
\end{equation}
что, по неравенству Грёнуолла, гарантирует $\|\bh_1(T)-\bh_2(T)\|\le e^{L_gT}\|\bx_1-\bx_2\|$. Это преобразует несертифицированное состязательное обучение уравнения~\eqref{eq:ps_robust} в задачу сертифицированной робастности. \emph{Подзадача 1 (Пробел~1) состоит в совместном решении уравнений~\eqref{eq:ps_robust}--\eqref{eq:ps_ode} с многомасштабной темпоральной динамикой и обучением с $O(1)$-памятью через сопряжённый метод. Решается Методом М1: CT-TGNN.}


\subsection{Этап 4: Распределённое обучение с приватностью}

В условиях распределённой динамики (Условие У3) граф $\calG$ разделён между $K$ участниками $\calD_1,\ldots,\calD_K$, которые не могут передавать сырые данные. Включение Требования Т4 (приватность) требует, чтобы алгоритм обучения удовлетворял $(\epsilon,\delta)$-дифференциальной приватности при устойчивости к доле $q$ византийских участников:
\begin{equation}
\min_\theta\;
\frac{1}{K}\sum_{k=1}^{K}\calL_k(\theta;\calD_k)
\quad\text{при условии}\quad
\Prob[\calM(\calD)\in S] \le e^{\epsilon}\Prob[\calM(\calD')\in S]+\delta,
\label{eq:ps_fed}
\end{equation}
где $\calM$ --- механизм обучения, и ограничение должно выполняться для всех соседних наборов данных $\calD,\calD'$. Кроме того, гетерогенные распределения участников требуют выравнивания через расстояние Вассерштейна $W_1(\calD_k,\calD_{k'})$, которое само должно вычисляться при ограничении приватности. \emph{Подзадача 2 (Пробел~2) состоит в решении уравнения~\eqref{eq:ps_fed} с византийско-устойчивой агрегацией и выравниванием оптимального транспорта при использовании сертификатов Липшица из Подзадачи~1. Решается Методом М2: FedLLM-API.}


\subsection{Этап 5: Многоуровневое представление и защита от забывания}

Включение Требования Т2 (точность) на темпоральных графах с иерархической структурой требует обучения представлений на нескольких семантических и временных масштабах --- уровнях компонентов, трассировок и узлов --- с предотвращением катастрофического забывания при эволюции топологии:
\begin{equation}
\min_\theta\;
\sum_{g\in\{s,r,n\}}\lambda_g\,\calL_g(\theta;\calG_g)
+ \lambda_{\mathrm{ewc}}\sum_j F_j(\theta_j-\theta_{j}^{*})^2
+ \lambda_{\mathrm{con}}\!\sum_{(g,g')}\mathrm{KL}(p_g\|p_{g'}),
\label{eq:ps_triple}
\end{equation}
где $F_j$ --- информация Фишера, штрафующая отклонение от ранее обученных параметров $\theta^*$, а член KL-дивергенции обеспечивает межуровневую согласованность. \emph{Подзадача 3 (Пробел~3) состоит в решении уравнения~\eqref{eq:ps_triple} со структурированной разреженностью и обучением с учётом квантования. Решается Методом М3: TripleE-TGNN.}


\subsection{Этап 6: Вычислительная эффективность}

Требование Т3 (эффективность) предписывает, чтобы сложность вывода составляла не более $\calO(L\log L)$ вместо $\calO(L^2)$ моделей на основе трансформеров. Это достигается переформулировкой последовательной обработки как селективной модели пространства состояний с входозависимыми переходами:
\begin{equation}
\bh_{k} = \bar{\bA}(u_k)\,\bh_{k-1} + \bar{\bB}(u_k)\,u_k,
\qquad
\bar{\bA}(u_k)=\exp(\bA(u_k)\Delta),
\label{eq:ps_ssm}
\end{equation}
свёрточное ядро которой вычисляется через быстрое преобразование Фурье. Это должно быть сочетано с прогрессивной состязательной дистилляцией робастности и PAC-байесовскими сертификатами обобщения. \emph{Подзадача 4 (Пробел~4) состоит в решении задачи эффективного вывода уравнения~\eqref{eq:ps_ssm} с формальными гарантиями робастности при отравлении на этапе обучения. Решается Методом М4: MambaShield.}


\subsection{Этап 7: Нестационарная адаптация с неопределённостью}

В условиях нестационарности (Условие У2) распределение $\calD_t$ изменяется во времени. Модель должна адаптироваться онлайн с количественной оценкой как эпистемической, так и алеаторной неопределённости и достижением сублинейного регрета:
\begin{equation}
R(T) = \max_{f^*}\sum_{t=1}^{T}\calL(f_{\theta_t},y_t) - \sum_{t=1}^{T}\calL(f^*,y_t) \le \calO(\sqrt{T\log|\calS|}),
\label{eq:ps_regret}
\end{equation}
где сравнение проводится с лучшей фиксированной стратегией ретроспективно. Предсказательная неопределённость декомпозируется как $U_{\mathrm{total}} = U_{\mathrm{epi}} + U_{\mathrm{ale}}$, при этом эпистемическая неопределённость управляет обнаружением дрейфа, а алеаторная неопределённость захватывает присущий шум.

\emph{Подзадача 5 (Пробел~5) касается прямой совместимости при смене протокольного режима и решается предметно-ориентированной фундаментальной моделью (CyberSecLLM). Подзадача 6 (Пробел~6) состоит в решении задачи онлайн-адаптации уравнения~\eqref{eq:ps_regret} с интегрированной декомпозицией эпистемической и алеаторной неопределённости и байесовскими апостериорными обновлениями. Решается Методами М5: Стохастический трансформер и М6: UC-HGP.}


\subsection{Полная единая задача}

Объединение всех этапов даёт единую многоцелевую задачу оптимизации настоящей диссертации:

\begin{equation}
\boxed{
\begin{aligned}
\min_\theta\;&
\underbrace{\E_{(\calG,y)\sim\calD_t}\!\Bigl[
  \max_{\|\bdelta\|\le\epsilon}
    \calL\!\bigl(f_\theta(\calG+\bdelta),y\bigr)
\Bigr]}_{\text{Состязат. робастность (Т1, У1)}}
+ \lambda_{\mathrm{acc}}\underbrace{\sum_g\calL_g(\theta;\calG_g)}_{\text{Многоуровн. точность (Т2)}}
\\[6pt]
&+ \lambda_{\mathrm{eff}}\underbrace{\mathrm{Complexity}(f_\theta)}_{\text{Эффективн. $\le\calO(L\log L)$ (Т3)}}
+ \lambda_{\mathrm{priv}}\underbrace{\mathrm{PrivCost}(\epsilon_{\mathrm{DP}},\delta)}_{\text{Дифф. приватность (Т4, У3)}}
\\[6pt]
&+ \lambda_{\mathrm{drift}}\underbrace{R(T)}_{\text{Нестац. регрет (У2)}}
+ \lambda_{\mathrm{ewc}}\underbrace{\sum_j F_j(\theta_j-\theta_j^*)^2}_{\text{Защита от забывания}}
+ \lambda_{\mathrm{con}}\underbrace{\sum_{i<j}\calL_{ij}(\theta_i,\theta_j)}_{\text{Связка методов (М7)}}
\end{aligned}
}
\label{eq:ps_unified}
\end{equation}

\noindent при ограничениях:
\begin{itemize}
\item Динамика ОДУ, ограниченная по Липшицу: $\|f_\theta(\bh_1,t)-f_\theta(\bh_2,t)\|\le L_g\|\bh_1-\bh_2\|$ (сертифиц. робастность)
\item Дифференциальная приватность: $\Prob[\calM(\calD)\in S]\le e^{\epsilon}\Prob[\calM(\calD')\in S]+\delta$ (приватность)
\item Византийская устойчивость: сходимость при доле $q<1/2$ повреждённых участников
\item Сублинейный регрет: $R(T)\le\calO(\sqrt{T\log|\calS|})$ (адаптация)
\end{itemize}

\noindent\textbf{Декомпозиция.} Данная единая задача неразрешима как монолитная оптимизация. Однако она допускает естественную декомпозицию на семь подзадач, по одной на каждый выявленный пробел, которые могут быть решены поочерёдной оптимизацией со связкой согласованности. Таблица~\ref{tab:decomposition} сопоставляет каждую подзадачу её пробелу, членам уравнения~\eqref{eq:ps_unified}, которые она адресует, и методу, её решающему.

\begin{table}[htbp]
\centering
\caption{Декомпозиция единой постановки задачи на семь подзадач.}
\label{tab:decomposition}
\small
\begin{tabular}{p{0.5cm} p{3.0cm} p{4.5cm} p{5.0cm}}
\toprule
\textbf{ПЗ} & \textbf{Метод} & \textbf{Адресуемые члены} & \textbf{Ключевое ограничение} \\
\midrule
1 & М1: CT-TGNN & Состязат. робастн. + динамика ОДУ & Сертификат Липшица: $e^{L_gT}\|\delta\bx\|$ \\
2 & М2: FedLLM-API & Стоимость приватн. + визант. агрегация & $(\epsilon,\delta)$-ДП + выравнивание ОТ \\
3 & М3: TripleE-TGNN & Многоуровн. точность + защита от забыв. & Штраф Фишера EWC + межуровневая KL \\
4 & М4: MambaShield & Эффективность + состязат. робастн. & $\calO(L\log L)$ + граница PAC-Байеса \\
5 & CyberSecLLM & Прямая совместимость (предметная) & Многозадачные блоки ODE-S6 \\
6 & М5/М6 & Нестац. регрет + неопределённость & Сублинейный регрет + декомп. эпи/але \\
7 & М7: Сертификация & Связка методов + теор.-игровая & Равновесие Штакельберга + согласованн. \\
\bottomrule
\end{tabular}
\end{table}

\emph{Подзадача 7 (Пробел~7) состоит в связке всех методов через член регуляризации согласованности $\sum_{i<j}\calL_{ij}$ в уравнении~\eqref{eq:ps_unified} и доказательстве сходимости поочерёдной оптимизации. Решается Методом М7: Сертификация и единой системой, разработанной в Разделе~\ref{sec:unified}.}

В оставшейся части данной главы каждый метод разрабатывается в порядке цепочки подзадач: М1 (Раздел~\ref{sec:ct_tgnn}), М2 (Раздел~\ref{sec:fedllm}), М3 (Раздел~\ref{sec:triple_e}), М4 (Раздел~\ref{sec:mamba}), М5/М6 (Раздел~\ref{sec:stochastic}), прямо совместимое обнаружение (Раздел~\ref{sec:pq_idps}) и М7 с единой системой (Раздел~\ref{sec:game}).



%%=====================================================================
\section{Метод М1: CT-TGNN --- Динамика темпоральных графовых нейронных сетей непрерывного времени}
\label{sec:ct_tgnn}
%%=====================================================================

Данный раздел решает Подзадачу~1 путём разработки графовой нейронной сети, представления узлов которой эволюционируют согласно нейронному обыкновенному дифференциальному уравнению, обеспечивая сертифицированную границу робастности из ограничения Липшица уравнения~\eqref{eq:ps_ode}. Метод адресует Пробел~1 (Состязательность $\times$ Робастность; Нестационарность $\times$ Робастность).

\end{equation}
где $\bh_v(t)\in\R^{d_h}$ --- вложение узла~$v$ в момент~$t$, $\calN_v(t)$ обозначает окрестность~$v$ при текущей топологии, $\mathbf{A}(t)$ кодирует структуру смежности, а $f_\theta$ --- нейронная сеть, параметризованная~$\theta$.

Для гетерогенных графов правая часть использует типоспецифичную передачу сообщений. Пусть $v$ принадлежит типу~$k$ и пусть $\calN_v^{(k')}(t)$ обозначает его соседей типа~$k'$. Динамика:
\begin{equation}
\frac{d\bh_v^{(k)}(t)}{dt}
= \sum_{k'=1}^{K}\sum_{u\in\calN_v^{(k')}(t)}
  \alpha_{uv}^{(k,k')}(t)\;\mathrm{MSG}_{k,k'}\!\bigl(
    \bh_u^{(k')}(t),\;\mathbf{e}_{uv}(t)\bigr),
\label{eq:ct_tgnn_hetero}
\end{equation}
где коэффициенты внимания~\cite{Velickovic2018}
\begin{equation}
\alpha_{uv}^{(k,k')}(t)
= \frac{
  \exp\!\bigl(\mathrm{LeakyReLU}\!\bigl(
    \mathbf{a}_{k,k'}^\top
    [\bW_k\bh_v^{(k)}(t)\|\bW_{k'}\bh_u^{(k')}(t)]
  \bigr)\bigr)
}{
  \sum_{w\in\calN_v^{(k')}(t)}
  \exp\!\bigl(\mathrm{LeakyReLU}\!\bigl(
    \mathbf{a}_{k,k'}^\top
    [\bW_k\bh_v^{(k)}(t)\|\bW_{k'}\bh_w^{(k')}(t)]
  \bigr)\bigr)
}
\label{eq:ct_tgnn_attn}
\end{equation}
адаптивно взвешивают важность соседей, $\bW_k,\bW_{k'}$ --- обучаемые проекции, а $\|$ обозначает конкатенацию.

\subsection{Многомасштабная темпоральная динамика}

Явления, разворачивающиеся на нескольких временных масштабах, от субсекундных переходных процессов до многомесячных тенденций, требуют параллельных ветвей ОДУ с различными постоянными времени. Для масштабов $s=1,\dots,S$ с постоянными $\tau_1<\tau_2<\cdots<\tau_S$ динамика:
\begin{equation}
\frac{d\bh_v^{(s)}(t)}{dt}
= \frac{1}{\tau_s}\;f_{\theta_s}\!\bigl(
  \bh_v^{(s)}(t),\;\{\bh_u^{(s)}(t)\}_{u\in\calN_v(t)},\;
  \mathbf{A}(t),\;t\bigr),
\label{eq:ct_tgnn_multi}
\end{equation}
и итоговое вложение агрегирует по масштабам через обученное внимание,
\begin{equation}
\bh_v(t)
= \sum_{s=1}^{S}\beta_s(t)\;\bh_v^{(s)}(t),
\qquad
\beta_s(t) = \mathrm{Softmax}_s\!\bigl(\mathbf{w}_s^\top\bh_v^{(s)}(t)\bigr),
\label{eq:ct_tgnn_scale_fuse}
\end{equation}
так что относительный вклад каждого временного масштаба адаптируется к локальной динамике.

\subsection{Метод сопряжённых чувствительностей для обучения с эффективным использованием памяти}

Обучение требует градиентов функции потерь $\calL=\calL(\bh(T))$ по~$\theta$ через решение ОДУ. Прямое обратное распространение через численный интегратор сохраняет все промежуточные состояния, что недопустимо для длинных горизонтов. Метод сопряжённых~\cite{Pontryagin1962,Chen2018} определяет сопряжённое состояние $\mathbf{a}(t)=\partial\calL/\partial\bh(t)$, удовлетворяющее обратному ОДУ
\begin{equation}
\frac{d\mathbf{a}(t)}{dt}
= -\mathbf{a}(t)^\top\frac{\partial f_\theta}{\partial\bh}(\bh(t),t),
\label{eq:adjoint_ode}
\end{equation}
интегрируемому от $\mathbf{a}(T)=\partial\calL/\partial\bh(T)$ до $\mathbf{a}(0)$. Градиенты по параметрам вычисляются из
\begin{equation}
\frac{\partial\calL}{\partial\theta}
= -\int_0^T \mathbf{a}(t)^\top
  \frac{\partial f_\theta}{\partial\theta}(\bh(t),t)\;dt.
\label{eq:adjoint_param}
\end{equation}
Затраты памяти пропорциональны размеру модели, а не длине траектории, что позволяет обучать на произвольно длинных темпоральных горизонтах.

\subsection{Темпоральная адаптивная пакетная нормализация}

Стандартная пакетная нормализация вычисляет статистики по фиксированному мини-батчу, что несовместимо с динамикой непрерывного времени, при которой распределение скрытых состояний непрерывно изменяется. Темпоральная адаптивная пакетная нормализация вычисляет скользящие статистики по временному окну $[t-\Delta,t]$,
\begin{equation}
\mathrm{TABN}\!\bigl(\bh(t)\bigr)
= \gamma(t)\;\frac{\bh(t)-\mu_{[t-\Delta,t]}}
  {\sqrt{\sigma^2_{[t-\Delta,t]}+\varepsilon}}
+ \beta(t),
\label{eq:tabn}
\end{equation}
где $\mu_{[t-\Delta,t]}$ и $\sigma^2_{[t-\Delta,t]}$ --- оконные среднее и дисперсия, а аффинные параметры $\gamma(t),\beta(t)$ сами эволюционируют через вспомогательные ОДУ $d\gamma/dt=g_\gamma(\bh(t))$, $d\beta/dt=g_\beta(\bh(t))$, управляемые обучаемыми сетями адаптации.

\subsection{Анализ сходимости и робастности}

\begin{theorem}[Сходимость CT-TGNN]
\label{thm:ct_tgnn_conv}
Предположим, что функция динамики $f_\theta$ является $L_f$-липшицевой по~$\bh$ и $L_\theta$-липшицевой по~$\theta$. При градиентном спуске со скоростью обучения $\eta<2/(L_f L_\theta)$ потери обучения сходятся к стационарной точке со скоростью $\calO(1/\sqrt{T})$, где $T$ обозначает число шагов градиента.
\end{theorem}

\noindent\textit{Эскиз доказательства.} Потери $\calL(\theta)$ вычисляются через карту решения ОДУ $\bh(T;\theta)$. По правилу цепочки и свойствам Липшица~$f_\theta$ градиент $\nabla_\theta\calL$ ограничен и липшицево-непрерывен, помещая задачу в стандартную схему гладкой невыпуклой оптимизации. Применение леммы убывания с указанной скоростью обучения даёт скорость $\calO(1/\sqrt{T})$ до $\epsilon$-стационарной точки. Полное доказательство приведено в Приложении~А.\qed

\begin{theorem}[Сертификат робастности по Липшицу]
\label{thm:ct_tgnn_rob}
Пусть $f_\theta$ является $L_g$-липшицевой по первому аргументу, а горизонт интегрирования равен~$T$. Для любых двух входов $\bx_1,\bx_2$ соответствующие вложения удовлетворяют
\begin{equation}
\|\bh_1(T)-\bh_2(T)\|
\le e^{L_g T}\;\|\bx_1-\bx_2\|.
\label{eq:lip_cert}
\end{equation}
Следовательно, если голова классификатора имеет константу Липшица~$L_c$, то сквозная модель является $L_c\,e^{L_g T}$-липшицевой, и предсказание сертифицировано постоянным в шаре радиуса $\rho = \Delta_{\min}/(L_c\,e^{L_g T})$, где $\Delta_{\min}$ --- запас между двумя наибольшими логитами.
\end{theorem}

\noindent\textit{Эскиз доказательства.} Пусть $\bm{\eta}(t)=\bh_1(t)-\bh_2(t)$. Тогда $d\bm{\eta}/dt = f_\theta(\bh_1,t)-f_\theta(\bh_2,t)$, откуда $d\|\bm{\eta}\|/dt\le L_g\|\bm{\eta}\|$ по свойству Липшица. Неравенство Грёнуолла даёт $\|\bm{\eta}(T)\|\le e^{L_g T}\|\bm{\eta}(0)\|$, а $\bm{\eta}(0)=\bx_1-\bx_2$. Композиция с головой классификатора завершает сертификат. Полный вывод со спектрально-нормовыми ограничениями, обеспечивающими границу Липшица при обучении, приведён в Приложении~А.\qed

\subsection{Вычислительная сложность}

Решение графового ОДУ адаптивным интегратором Рунге--Кутты при допуске~$\varepsilon_{\mathrm{ode}}$ требует $\calO(\varepsilon_{\mathrm{ode}}^{-1/5})$ вычислений функции. Каждое вычисление выполняет передачу сообщений за $\calO(|\calV|\,d\,d_h)$ при ограниченной степени~$d$ и размерности вложения~$d_h$. Общая сложность обучения за эпоху, таким образом:
\begin{equation}
\calO\!\bigl(|\calV|\,d\,d_h\;\varepsilon_{\mathrm{ode}}^{-1/5}\;T_{\mathrm{train}}\bigr),
\label{eq:ct_tgnn_complex}
\end{equation}
что близко к линейной по размеру графа при $|\calE|=\calO(|\calV|)$.

\begin{algorithm}[htbp]
\caption{Обучение CT-TGNN с методом сопряжённых чувствительностей}\label{alg:ct_tgnn}
\KwInput{Последовательность темпоральных графов $\{\calG(t_i)\}_{i=1}^N$, метки $\{y_i\}$, скорость обучения $\eta$, горизонт интегрирования $T$, допуск $\varepsilon_{\mathrm{ode}}$}
\KwOutput{Обученные параметры $\theta^*$}
Инициализировать параметры $\theta_0$ и многомасштабные постоянные времени $\tau_1<\cdots<\tau_S$\;
\For{эпоха $= 1,2,\ldots,E$}{
  \For{каждый мини-батч $\{(\calG(t_j),y_j)\}$}{
    Установить начальные вложения $\bh_v(0) \leftarrow \mathrm{Embed}(\bx_v(t_j))$ для всех $v\in\calV$\;
    \For{масштаб $s=1,\ldots,S$}{
      Решить $d\bh_v^{(s)}/dt = \tau_s^{-1}f_{\theta_s}(\bh_v^{(s)},\calN_v,\bA,t)$ от $0$ до $T$ адаптивным РК с допуском $\varepsilon_{\mathrm{ode}}$\;
    }
    Вычислить объединённые вложения $\bh_v(T) = \sum_s \beta_s(T)\,\bh_v^{(s)}(T)$\;
    Вычислить потери классификации $\calL = \calL_{\mathrm{CE}}(f_{\mathrm{head}}(\bh(T)),\,y)$\;
    Установить терминальное условие сопряжённого $\mathbf{a}(T) = \partial\calL/\partial\bh(T)$\;
    Решить сопряжённое ОДУ назад: $d\mathbf{a}/dt = -\mathbf{a}^\top(\partial f_\theta/\partial\bh)$ от $T$ до $0$\;
    Накопить градиенты по параметрам $\partial\calL/\partial\theta = -\int_0^T \mathbf{a}^\top(\partial f_\theta/\partial\theta)\,dt$\;
    Обновить $\theta \leftarrow \theta - \eta\,\partial\calL/\partial\theta$\;
  }
}
Вычислить сертификат Липшица $\rho = \Delta_{\min}/(L_c\,e^{L_g T})$ для каждого тестового входа\;
\Return $\theta^*,\;\{\rho_i\}$
\end{algorithm}


%%=====================================================================
\section{Метод М2: FedLLM-API --- Федеративная оптимизация графов с приватностью}
\label{sec:fedllm}
%%=====================================================================

Данный раздел решает Подзадачу~2 путём разработки алгоритма федеративного обучения, обеспечивающего ограничение приватности и византийскую устойчивость уравнения~\eqref{eq:ps_fed} при использовании сертификатов Липшица, полученных М1. Метод адресует Пробел~2 (Распределённость $\times$ Робастность; Распределённость $\times$ Приватность; Распределённость $\times$ Эффективность).

\subsection{Византийско-устойчивая агрегация}

Рассмотрим $K$ участников с локальными наборами данных $\calD_1,\dots,\calD_K$, из которых не более доли~$q$ являются византийскими. На раунде коммуникации~$t$ сервер рассылает глобальную модель~$\bw_t$, и каждый честный клиент~$k$ вычисляет локальное обновление через градиентный спуск на своих данных, $\bw_{t+1}^{(k)}=\bw_t - \eta\nabla_{\bw}\calL_k(\bw_t;\calD_k)$. Злонамеренные клиенты отправляют произвольные векторы. Сервер агрегирует через покоординатное усечённое среднее~\cite{Yin2018,Blanchard2017,Li2020fedprox}: для каждого измерения параметра~$j$ значения клиентов сортируются и $b=\lfloor Kq\rfloor$ наибольших и наименьших отбрасываются, давая
\begin{equation}
\bar{w}_{t+1,j}
= \frac{1}{K-2b}\sum_{i=b+1}^{K-b}w_{t+1,j}^{(i)},
\label{eq:fed_trim}
\end{equation}
где $w_{t+1,j}^{(i)}$ --- $i$-я порядковая статистика.

\begin{theorem}[Византийско-устойчивая сходимость]
\label{thm:fed_conv}
Предположим, что глобальные потери $F(\bw)=\frac{1}{K}\sum_{k=1}^K F_k(\bw)$ являются $\mu$-сильно выпуклыми и $L$-гладкими, а доля византийских клиентов удовлетворяет $f<q<1/2$. При скорости обучения $\eta\le 1/L$ покоординатное усечённое среднее достигает
\begin{equation}
\E\!\bigl[\|\bw_T-\bw^*\|^2\bigr]
\le \Bigl(1-\frac{\mu\eta}{2}\Bigr)^{\!T}\|\bw_0-\bw^*\|^2
  + \frac{2\eta}{\mu(K{-}2b)}
    \Bigl(\sum_{k=1}^{K-2b}\sigma_k^2 + C_{\mathrm{byz}}\Bigr),
\label{eq:fed_conv_bound}
\end{equation}
где $\sigma_k^2$ --- дисперсия градиента честного клиента~$k$, а $C_{\mathrm{byz}}$ ограничивает влияние византийских участников.
\end{theorem}

\noindent\textit{Эскиз доказательства.} Ключевой шаг --- показать, что усечённое среднее градиентных векторов честных и византийских участников близко к истинному среднему честных градиентов. Пусть $\bar{\mathbf{g}}_H$ --- среднее честных градиентов и $\bar{\mathbf{g}}_{\mathrm{TM}}$ --- выход усечённого среднего. Поскольку не более $b$ из $2b$ отброшенных записей могут быть честными, ошибка $\|\bar{\mathbf{g}}_{\mathrm{TM}}-\bar{\mathbf{g}}_H\|$ ограничена порядковыми статистиками норм честных градиентов через неравенства концентрации. Подстановка в стандартную рекурсию убывания для сильно выпуклых функций даёт указанную границу. Полные детали приведены в Приложении~А.\qed

\subsection{Дифференциальная приватность через калиброванный шум}

Каждый клиент обрезает норму градиента до порога~$C$ и добавляет гауссовский шум~\cite{Abadi2016,Dwork2006,Dwork2014} перед передачей,
\begin{equation}
\bar{\mathbf{g}}_k = \mathbf{g}_k\Big/\max\!\Bigl(1,\frac{\|\mathbf{g}_k\|}
{C}\Bigr),
\qquad
\tilde{\mathbf{g}}_k = \bar{\mathbf{g}}_k + \calN(\mathbf{0},\sigma^2 C^2\mathbf{I}),
\label{eq:fed_dp_noise}
\end{equation}
где масштаб шума~$\sigma$ определяется целевыми параметрами приватности $(\epsilon,\delta)$. Усиление приватности подвыборкой с вероятностью батча~$q_b$ и композиция за $T$ раундов через моменты-бухгалтер~\cite{Mironov2017,Balle2018} даёт общую гарантию приватности
\begin{equation}
\epsilon_{\mathrm{total}}
= \min_{\alpha>1}\;
  \frac{1}{\alpha-1}\log\frac{1}{\delta}
  + \frac{T\,q_b^2\,\alpha}{2\sigma^2}.
\label{eq:fed_privacy}
\end{equation}

\subsection{Оптимальный транспорт для выравнивания распределений}

Гетерогенные распределения клиентов ухудшают сходимость. Оптимальный транспорт выравнивает их путём вычисления Вассерштейнова барицентра. Для дискретных исходных выборок $\{\bx_i\}_{i=1}^{n_s}$ и целевых выборок $\{\by_j\}_{j=1}^{n_t}$ дискретный транспортный план решает
\begin{equation}
\mathbf{T}^* = \arg\min_{\mathbf{T}\in\R_+^{n_s\times n_t}}
\langle\mathbf{C},\mathbf{T}\rangle_F
\quad\text{при}\quad
\mathbf{T}\mathbf{1}_{n_t}=\mathbf{a},\;
\mathbf{T}^\top\mathbf{1}_{n_s}=\mathbf{b},
\label{eq:fed_ot}
\end{equation}
где $C_{ij}=c(\bx_i,\by_j)$ --- матрица стоимости, а $\mathbf{a},\mathbf{b}$ лежат на вероятностных симплексах. Алгоритм Синкхорна~\cite{Cuturi2013,Chizat2020,Villani2009} вычисляет энтропийно регуляризованное решение через поочерёдное масштабирование,
\begin{equation}
\mathbf{u}^{(\ell+1)} = \mathbf{a}\oslash(\mathbf{K}\mathbf{v}^{(\ell)}),
\qquad
\mathbf{v}^{(\ell+1)} = \mathbf{b}\oslash(\mathbf{K}^\top\mathbf{u}^{(\ell+1)}),
\label{eq:sinkhorn}
\end{equation}
где $\mathbf{K}=\exp(-\mathbf{C}/\lambda)$ --- энтропийное ядро, $\oslash$ обозначает поэлементное деление, а $\lambda>0$ управляет регуляризацией. После сходимости транспортный план $\mathbf{T}^*=\mathrm{diag}(\mathbf{u})\mathbf{K}\,\mathrm{diag}(\mathbf{v})$.

\subsection{Дифференциально приватный оптимальный транспорт}

Для защиты расположения выборок гауссовский шум добавляется к эмпирическому распределению перед вычислением транспорта,
\begin{equation}
\tilde{\mu}_s = \frac{1}{n_s}\sum_{i=1}^{n_s}\delta_{\bx_i+\bm{\xi}_i},
\qquad \bm{\xi}_i\sim\calN(\mathbf{0},\sigma_{\mathrm{DP}}^2\mathbf{I}),
\label{eq:fed_dp_ot_dist}
\end{equation}
а шум Лапласа добавляется к результирующему транспортному плану,
\begin{equation}
\tilde{\mathbf{T}} = \mathbf{T}^*
+ \mathrm{Lap}\!\Bigl(\frac{\Delta_T}{\epsilon_{\mathrm{OT}}}\Bigr),
\label{eq:fed_dp_ot_plan}
\end{equation}
где чувствительность $\Delta_T=\max_{\calD,\calD'}\|\mathbf{T}^*(\calD)-\mathbf{T}^*(\calD')\|_1$ ограничивает максимальное изменение при замене одной записи, а $\epsilon_{\mathrm{OT}}$ --- бюджет приватности, выделенный на транспорт.

\subsection{Интеграция большой языковой модели для обобщения zero-shot}

Семантическое рассуждение о новых категориях достигается построением промпта на естественном языке из структурных признаков наблюдения и запросом к большой языковой модели для получения вероятности $p_{\mathrm{LLM}}(\mathrm{class}\mid\mathbf{r})$. Итоговое предсказание --- ансамбль
\begin{equation}
p_{\mathrm{final}}
= \lambda_{\mathrm{LLM}}\;p_{\mathrm{LLM}}(\mathrm{class}\mid\mathbf{r})
  + (1-\lambda_{\mathrm{LLM}})\;p_{\mathrm{GNN}}(\mathrm{class}\mid\calG(\mathbf{r})),
\label{eq:fed_llm}
\end{equation}
где $\calG(\mathbf{r})$ --- граф, построенный из наблюдения, а $\lambda_{\mathrm{LLM}}$ балансирует два источника.

\begin{algorithm}[htbp]
\caption{FedLLM-API: Византийско-устойчивая федеративная оптимизация с ДП-ОТ}\label{alg:fedllm}
\KwInput{Клиенты $\calD_1,\ldots,\calD_K$, доля византийских $q$, порог обрезки $C$, бюджет приватности $(\epsilon,\delta)$, регуляризация Синкхорна $\lambda$, раунды коммуникации $T$}
\KwOutput{Глобальная модель $\bw_T$}
Инициализировать глобальную модель $\bw_0$; вычислить масштаб шума $\sigma = C\sqrt{2\ln(1.25/\delta)}/\epsilon$\;
\For{раунд $t = 0,1,\ldots,T{-}1$}{
  Разослать $\bw_t$ всем клиентам\;
  \For{каждый клиент $k = 1,\ldots,K$ \textup{(параллельно)}}{
    Вычислить локальный градиент $\mathbf{g}_k = \nabla_{\bw}\calL_k(\bw_t;\calD_k)$\;
    Обрезать: $\bar{\mathbf{g}}_k \leftarrow \mathbf{g}_k / \max(1,\|\mathbf{g}_k\|/C)$\;
    Добавить шум: $\tilde{\mathbf{g}}_k \leftarrow \bar{\mathbf{g}}_k + \calN(\mathbf{0},\sigma^2 C^2\mathbf{I})$\;
    Отправить $\tilde{\mathbf{g}}_k$ серверу\;
  }
  Установить $b \leftarrow \lfloor Kq\rfloor$\;
  \For{каждое измерение $j$}{
    Отсортировать $\{\tilde{g}_{k,j}\}_{k=1}^K$; вычислить усечённое среднее $\bar{g}_{t,j} = \frac{1}{K-2b}\sum_{i=b+1}^{K-b}\tilde{g}_{j}^{(i)}$\;
  }
  \If{обнаружена гетерогенность по оценке $W_1$}{
    Вычислить зашумлённые исходные распределения $\tilde{\mu}_k$ добавлением $\calN(\mathbf{0},\sigma_{\mathrm{DP}}^2\mathbf{I})$ к выборкам\;
    Выполнить итерации Синкхорна для вычисления регуляризованного транспортного плана $\mathbf{T}^*$\;
    Добавить шум Лапласа: $\tilde{\mathbf{T}} \leftarrow \mathbf{T}^* + \mathrm{Lap}(\Delta_T/\epsilon_{\mathrm{OT}})$\;
    Применить коррекцию выравнивания к $\bar{\mathbf{g}}_t$ с использованием $\tilde{\mathbf{T}}$\;
  }
  Обновить глобальную модель: $\bw_{t+1} \leftarrow \bw_t - \eta\,\bar{\mathbf{g}}_t$\;
}
\Return $\bw_T$
\end{algorithm}


%%=====================================================================
\section{Метод М3: TripleE-TGNN --- Многоуровневые темпоральные вложения графов}
\label{sec:triple_e}
%%=====================================================================

Данный раздел решает Подзадачу~3 путём разработки многоуровневых вложений графов, адресующих члены точности и защиты от забывания в уравнении~\eqref{eq:ps_triple}. Метод адресует Пробел~3 (Нестационарность $\times$ Точность; Нестационарность $\times$ Эффективность).

\subsection{Иерархическое представление графов}

Многие системы обладают естественной иерархией: грубый уровень функциональных компонентов, промежуточный уровень операционных трассировок, охватывающих компоненты, и тонкий уровень узлов инфраструктуры. Эта структура захватывается тремя связанными темпоральными графами: графом компонентов $\calG_s(t)=(\calV_s,\calE_s(t),\mathbf{X}_s(t))$, графом трассировок $\calG_r(t)=(\calV_r,\calE_r(t),\mathbf{X}_r(t))$ и графом узлов $\calG_n(t)=(\calV_n,\calE_n(t),\mathbf{X}_n(t))$. Двудольные графы связи $\calB_{sr},\calB_{sn},\calB_{rn}$ соединяют сущности между уровнями: компонент развёрнут на узле, трассировка проходит через компонент, трассировка выполняется на узле.

\subsection{Многоуровневое обучение вложений}

Каждая гранулярность $g\in\{s,r,n\}$ поддерживает выделенный кодировщик,
\begin{equation}
\bh_v^{(g)}(t)
= \mathrm{ENC}_g\!\bigl(
  \bx_v^{(g)}(t),\;
  \{\bh_u^{(g)}(t{-}1)\}_{u\in\calN_v^{(g)}(t)},\;
  \Delta t
\bigr),
\label{eq:triple_enc}
\end{equation}
где $\mathrm{ENC}_g$ --- графовая нейронная сеть, специфичная для данной гранулярности. Межгранулярная передача сообщений интегрирует информацию между уровнями абстракции через двудольное внимание,
\begin{equation}
\mathbf{m}_v^{(g\to g')}(t)
= \sum_{u\in\calN_v^{(g,g')}}
  \alpha_{uv}^{(g,g')}(t)\;\bW_{g\to g'}\bh_u^{(g)}(t),
\label{eq:triple_cross_msg}
\end{equation}
с коэффициентами внимания
\begin{equation}
\alpha_{uv}^{(g,g')}(t)
= \mathrm{Softmax}_u\!\Biggl(
  \frac{(\mathbf{Q}_{g'}\bh_v^{(g')}(t))^\top
        (\mathbf{K}_g\bh_u^{(g)}(t))}
       {\sqrt{d_k}}
\Biggr),
\label{eq:triple_cross_attn}
\end{equation}
где $\mathbf{Q}_{g'},\mathbf{K}_g$ --- проекции запроса и ключа, а $\calN_v^{(g,g')}$ обозначает соседей в графе связи. Итоговое вложение объединяет внутригранулярную и межгранулярную информацию через нормализацию слоя,
\begin{equation}
\bh_v^{(g)}(t)
\leftarrow
\mathrm{LayerNorm}\!\Bigl(
  \bh_v^{(g)}(t)
  + \sum_{g'\neq g}\mathbf{m}_v^{(g'\to g)}(t)
\Bigr).
\label{eq:triple_fuse}
\end{equation}

\subsection{Двойное темпорально-пространственное внимание}

Внутри каждой гранулярности темпоральное внимание агрегирует исторические состояния по окну ширины~$w$,
\begin{equation}
\tilde{\bh}_v^{(g)}(t)
= \sum_{\tau=t-w}^{t-1}
  \alpha_\tau^{\mathrm{temp}}\;\bh_v^{(g)}(\tau),
\qquad
\alpha_\tau^{\mathrm{temp}}
= \mathrm{Softmax}_\tau\!\bigl(
  \mathbf{w}_{\mathrm{temp}}^\top\bh_v^{(g)}(\tau)\bigr),
\label{eq:triple_temp_attn}
\end{equation}
а пространственное внимание агрегирует текущую окрестность,
\begin{equation}
\hat{\bh}_v^{(g)}(t)
= \sum_{u\in\calN_v^{(g)}(t)}
  \alpha_u^{\mathrm{spat}}\;\bh_u^{(g)}(t),
\qquad
\alpha_u^{\mathrm{spat}}
= \mathrm{Softmax}_u\!\bigl(
  \mathbf{w}_{\mathrm{spat}}^\top
  [\bh_v^{(g)}(t)\|\bh_u^{(g)}(t)]\bigr).
\label{eq:triple_spat_attn}
\end{equation}
Комбинированное представление:
\begin{equation}
\bh_v^{(g)}(t)
= \mathrm{MLP}\!\bigl(
  [\tilde{\bh}_v^{(g)}(t)\|
   \hat{\bh}_v^{(g)}(t)\|
   \bh_v^{(g)}(t)]
\bigr).
\label{eq:triple_combined}
\end{equation}

\subsection{Упругая консолидация весов для эволюции топологии}

При изменении топологии графа, например при развёртывании или реконфигурации инфраструктуры, модель, переобученная на новой топологии, может перезаписать ранее обученные представления. Упругая консолидация весов~\cite{Kirkpatrick2017,Zenke2017} предотвращает это катастрофическое забывание путём штрафования изменений параметров, важных для предыдущих задач. После обучения на конфигурации топологии~$\calT_i$ с оптимальными параметрами~$\theta_i^*$, потери для последующей конфигурации~$\calT_{i+1}$ включают член регуляризации
\begin{equation}
\calL_{\mathrm{EWC}}(\theta)
= \calL_{\calT_{i+1}}(\theta)
  + \lambda_{\mathrm{ewc}}\sum_j F_j\,(\theta_j - \theta_{i,j}^*)^2,
\label{eq:ewc}
\end{equation}
где диагональная информация Фишера $F_j = \E[(\partial\log p(y|\bx,\theta_i^*)/\partial\theta_j)^2]$ оценивает важность параметра~$j$ для предыдущей конфигурации, а $\lambda_{\mathrm{ewc}}$ управляет силой консолидации.

\subsection{Совместная многоуровневая целевая функция классификации}

Модель выполняет одновременную классификацию на всех трёх гранулярностях через многозадачные потери,
\begin{equation}
\calL_{\mathrm{joint}}
= \lambda_s\calL_s + \lambda_r\calL_r + \lambda_n\calL_n
  + \lambda_{\mathrm{con}}\calL_{\mathrm{consist}},
\label{eq:triple_joint}
\end{equation}
где $\calL_g$ --- кросс-энтропийные потери на гранулярности~$g$, $\lambda_g$ --- балансирующие веса, а потери согласованности
\begin{equation}
\calL_{\mathrm{consist}}
= \sum_{(v_s,v_r)\in\calB_{sr}}
  \mathrm{KL}\!\bigl(p_s(y|v_s)\|p_r(y|v_r)\bigr)
+ \sum_{(v_s,v_n)\in\calB_{sn}}
  \mathrm{KL}\!\bigl(p_s(y|v_s)\|p_n(y|v_n)\bigr)
\label{eq:triple_consist}
\end{equation}
обеспечивают согласованность предсказаний связанных сущностей между уровнями абстракции.

\begin{algorithm}[htbp]
\caption{Многоуровневое обучение TripleE-TGNN}\label{alg:triple_e}
\KwInput{Иерархические графы $\{\calG_s,\calG_r,\calG_n\}$ со связями $\calB_{sr},\calB_{sn},\calB_{rn}$, метки на каждой гранулярности, сила EWC $\lambda_{\mathrm{ewc}}$}
\KwOutput{Обученные кодировщики $\mathrm{ENC}_s,\mathrm{ENC}_r,\mathrm{ENC}_n$ и параметры межуровневого внимания}
Инициализировать все параметры кодировщиков и внимания\;
\For{каждая конфигурация топологии $\calT_i$}{
  \For{эпоха $= 1,\ldots,E$}{
    \For{каждый мини-батч}{
      \For{гранулярность $g \in \{s,r,n\}$}{
        Вычислить внутригранулярные вложения $\bh_v^{(g)}(t) \leftarrow \mathrm{ENC}_g(\bx_v^{(g)},\calN_v^{(g)},\Delta t)$\;
        Вычислить темпоральное внимание $\tilde{\bh}_v^{(g)}$ и пространственное внимание $\hat{\bh}_v^{(g)}$\;
        Объединить: $\bh_v^{(g)} \leftarrow \mathrm{MLP}([\tilde{\bh}_v^{(g)}\|\hat{\bh}_v^{(g)}\|\bh_v^{(g)}])$\;
      }
      \For{каждая пара $(g,g')$ с $g\neq g'$}{
        Вычислить межгранулярные сообщения $\mathbf{m}_v^{(g\to g')}$ через двудольное внимание\;
        Обновить: $\bh_v^{(g')} \leftarrow \mathrm{LayerNorm}(\bh_v^{(g')} + \mathbf{m}_v^{(g\to g')})$\;
      }
      Вычислить $\calL_{\mathrm{joint}} = \lambda_s\calL_s + \lambda_r\calL_r + \lambda_n\calL_n + \lambda_{\mathrm{con}}\calL_{\mathrm{consist}}$\;
      \If{$i > 1$}{
        Добавить штраф EWC: $\calL \leftarrow \calL_{\mathrm{joint}} + \lambda_{\mathrm{ewc}}\sum_j F_j(\theta_j-\theta_{i-1,j}^*)^2$\;
      }
      Обратное распространение и обновление всех параметров\;
    }
  }
  Сохранить диагональ Фишера $\{F_j\}$ и оптимальные $\theta_i^*$ для EWC\;
}
\Return обученные кодировщики и параметры внимания
\end{algorithm}



%%=====================================================================
\section{Метод М4: MambaShield --- Вывод на основе пространства состояний с устойчивостью к отравлению}
\label{sec:mamba}
%%=====================================================================

Данный раздел решает Подзадачу~4 путём разработки эффективной последовательной модели, удовлетворяющей ограничению вычислительной сложности уравнения~\eqref{eq:ps_ssm} при сохранении робастности при отравлении на этапе обучения через прогрессивную дистилляцию и PAC-байесовскую сертификацию. Метод адресует Пробел~4 (Состязательность $\times$ Эффективность; Нестационарность $\times$ Эффективность).

\subsection{Формулировка селективного пространства состояний}

Селективная модель пространства состояний~\cite{Gu2022,Gu2024} расширяет классическую линейную систему через входозависимые переходы. Для входной последовательности $u(t)\in\R^{d_{\mathrm{in}}}$ динамика непрерывного времени:
\begin{equation}
\frac{d\bh(t)}{dt} = \bA(u(t))\,\bh(t) + \bB(u(t))\,u(t),
\qquad
y(t) = \bC\,\bh(t) + \bD\,u(t),
\label{eq:mamba_cts}
\end{equation}
где $\bA(u)=f_A(u;\theta_A)\in\R^{d_s\times d_s}$ и $\bB(u)=f_B(u;\theta_B)\in\R^{d_s\times d_{\mathrm{in}}}$ вычисляются обученными проекциями. Дискретизация через удержание нулевого порядка с шагом~$\Delta$ даёт
\begin{equation}
\bh_k = \bar{\bA}_k\,\bh_{k-1} + \bar{\bB}_k\,u_k,
\qquad
\bar{\bA}_k = \exp(\bA(u_k)\Delta),
\quad
\bar{\bB}_k = \bA(u_k)^{-1}(\bar{\bA}_k - \mathbf{I})\bB(u_k),
\label{eq:mamba_disc}
\end{equation}
и свёрточное ядро $\bar{\mathbf{K}}=(\bC\bar{\bA}^i\bar{\bB})_{i=0}^{L-1}$ вычисляется за время $\calO(L\log L)$ через быстрое преобразование Фурье.

\subsection{Прогрессивная состязательная дистилляция робастности}

Ансамбль из $M$ моделей-учителей $\{f_{\theta_m}\}_{m=1}^M$, каждая обученная на непересекающемся подмножестве данных, обеспечивает робастность через дистилляцию знаний~\cite{Hinton2015,Xie2019}. Модель-ученик~$f_\phi$ обучается у учителей через потери
\begin{equation}
\calL_{\mathrm{distill}}
= \sum_{m=1}^{M}w_m\;\mathrm{KL}\!\bigl(p_\phi(y|\bx)\|p_{\theta_m}(y|\bx)\bigr),
\label{eq:mamba_distill}
\end{equation}
где веса учителей~$w_m$ отражают производительность на валидации. Учебный план $\rho(t)=\min(\rho_0+\alpha t,\rho_{\max})$ постепенно увеличивает долю повреждённых образцов в каждом мини-батче, а общие потери обучения интерполируют между потерями задачи и потерями дистилляции,
\begin{equation}
\calL_{\mathrm{total}}
= (1-\lambda(t))\,\calL_{\mathrm{task}}
+ \lambda(t)\,\calL_{\mathrm{distill}},
\qquad
\lambda(t) = \min(\lambda_0+\beta t,\;1),
\label{eq:mamba_curriculum}
\end{equation}
так что влияние учителей возрастает с увеличением интенсивности отравления.

\subsection{PAC-байесовская граница обобщения}

\begin{theorem}[PAC-байесовская граница для селективных моделей пространства состояний~\cite{McAllester1999,Catoni2007}]
\label{thm:pac_bayes}
Для любого $\delta\in(0,1)$, с вероятностью не менее $1-\delta$ по выбору обучающего множества~$S$ размера~$n$, для всех апостериорных распределений~$Q$ по параметрам,
\begin{equation}
\E_{\theta\sim Q}[R(\theta)]
\le \E_{\theta\sim Q}[\hat{R}_n(\theta)]
+ \sqrt{\frac{\mathrm{KL}(Q\|P)+\log(2\sqrt{n}/\delta)}{2n}},
\label{eq:pac_bayes}
\end{equation}
где $R(\theta)$ --- истинный риск, $\hat{R}_n(\theta)$ --- эмпирический риск, а $P$ --- не зависящее от данных априорное распределение.
\end{theorem}

\noindent\textit{Эскиз доказательства.} Доказательство следует стандартной PAC-байесовской схеме Катони, применяя замену меры с $P$ на~$Q$ через вариационное представление KL-дивергенции Донскера--Варадхана в сочетании с аргументом производящей функции моментов для эмпирического процесса. Граница справедлива равномерно по всем апостериорным~$Q$, включая сосредоточенные на единственном параметрическом векторе, и поэтому невакуумна, когда член KL мал относительно $n$. При гауссовском априорном $P=\calN(\mathbf{0},\sigma_P^2\mathbf{I})$ и апостериорном $Q=\calN(\bm{\mu}_Q,\sigma_Q^2\mathbf{I})$ член KL имеет замкнутую форму
\begin{equation}
\mathrm{KL}(Q\|P)
= \frac{d}{2}\Bigl(
  \frac{\sigma_Q^2}{\sigma_P^2}
  + \frac{\|\bm{\mu}_Q\|^2}{\sigma_P^2}
  - 1
  - \log\frac{\sigma_Q^2}{\sigma_P^2}
\Bigr),
\label{eq:pac_kl_gauss}
\end{equation}
который минимизируется сохранением апостериорного близко к началу координат и его дисперсии близко к дисперсии априорного. Полные детали приведены в Приложении~А.\qed

\begin{algorithm}[htbp]
\caption{MambaShield: Обучение с прогрессивной состязательной дистилляцией робастности}\label{alg:mamba}
\KwInput{Обучающие данные $\calD$, ансамбль учителей $\{f_{\theta_m}\}_{m=1}^M$, учебный план $\rho_0,\alpha,\rho_{\max}$, априорное PAC-Байеса $P$}
\KwOutput{Робастная модель-ученик $f_\phi$ с PAC-байесовским сертификатом}
Инициализировать параметры ученика $\phi_0$\;
Разделить $\calD$ на $M$ непересекающихся подмножеств; обучить каждого учителя $f_{\theta_m}$ на его подмножестве\;
Вычислить веса учителей $w_m \propto \mathrm{ValAcc}(f_{\theta_m})$\;
\For{итерация $t = 1,\ldots,T_{\mathrm{train}}$}{
  Установить долю отравления $\rho(t) = \min(\rho_0 + \alpha t,\;\rho_{\max})$\;
  Установить вес дистилляции $\lambda(t) = \min(\lambda_0 + \beta t,\;1)$\;
  Выбрать мини-батч; повредить долю $\rho(t)$ меток\;
  \For{каждый образец $(\bx,y)$}{
    Дискретизовать SSM: $\bar{\bA}_k = \exp(\bA(u_k)\Delta)$, $\bar{\bB}_k = \bA^{-1}(\bar{\bA}_k - \mathbf{I})\bB(u_k)$\;
    Выполнить селективное сканирование: $\bh_k = \bar{\bA}_k\bh_{k-1} + \bar{\bB}_k u_k$ для $k=1,\ldots,L$\;
    Вычислить предсказание ученика $p_\phi(y|\bx)$\;
  }
  Вычислить $\calL_{\mathrm{task}} = \mathrm{CE}(p_\phi,y)$\;
  Вычислить $\calL_{\mathrm{distill}} = \sum_m w_m\,\mathrm{KL}(p_\phi\|p_{\theta_m})$\;
  Итого: $\calL = (1-\lambda(t))\calL_{\mathrm{task}} + \lambda(t)\calL_{\mathrm{distill}} + \gamma\,\mathrm{KL}(Q_\phi\|P)$\;
  Обновить $\phi \leftarrow \phi - \eta\nabla_\phi\calL$\;
}
Вычислить границу PAC-Байеса: $R \le \hat{R}_n + \sqrt{(\mathrm{KL}(Q\|P)+\log(2\sqrt{n}/\delta))/(2n)}$\;
\Return $f_\phi$ и сертификат обобщения
\end{algorithm}


%%=====================================================================
\section{Методы М5 и М6: Вывод с калиброванной неопределённостью на графах}
\label{sec:stochastic}
%%=====================================================================

Данный раздел решает Подзадачу~6 путём разработки двух взаимодополняющих методов, обеспечивающих сублинейный регрет и декомпозицию неопределённости, требуемые уравнением~\eqref{eq:ps_regret}. Метод М5 (Стохастический трансформер) обеспечивает вариационное байесовское внимание, а Метод М6 (UC-HGP) --- неопределённость на иерархических гауссовских процессах на нескольких временных масштабах. Вместе они адресуют Пробел~6 (Состязательность $\times$ Точность; Нестационарность $\times$ Точность).

\subsection{Вариационное байесовское внимание}

Для входной последовательности $\mathbf{X}=[\bx_1,\dots,\bx_n]\in\R^{n\times d}$ вариационные распределения~\cite{Blundell2015} размещаются на матрицах проекций запроса, ключа и значения,
\begin{equation}
\begin{aligned}
q_{\phi_Q}(\bW_Q) = \calN(\bm{\mu}_Q,\mathrm{diag}(\bm{\sigma}_Q^2)),
\quad
q_{\phi_K}(\bW_K) = \calN(\bm{\mu}_K,\mathrm{diag}(\bm{\sigma}_K^2)), \\
\quad
q_{\phi_V}(\bW_V) = \calN(\bm{\mu}_V,\mathrm{diag}(\bm{\sigma}_V^2)).
\end{aligned}
\label{eq:stoch_var_attn}
\end{equation}
Выборки производятся через трюк репараметризации, $\bW_Q=\bm{\mu}_Q+\bm{\sigma}_Q\odot\bm{\epsilon}_Q$ с $\bm{\epsilon}_Q\sim\calN(\mathbf{0},\mathbf{I})$, так что коэффициенты внимания $A_{ij}=(\bx_i\bW_Q)^\top(\bx_j\bW_K)/\sqrt{d_k}$ становятся случайными величинами, неопределённость которых распространяется на выход. Оценка Монте-Карло с $M$ выборками аппроксимирует ожидаемое внимание,
\begin{equation}
\E_{q_\phi}[\mathrm{Attention}(\mathbf{X})]
\approx \frac{1}{M}\sum_{m=1}^{M}
  \mathrm{Attention}\!\bigl(\mathbf{X};\bW_Q^{(m)},\bW_K^{(m)},\bW_V^{(m)}\bigr).
\label{eq:stoch_mc_attn}
\end{equation}

\subsection{Декомпозиция эпистемической и алеаторной неопределённости}

Полная предсказательная неопределённость декомпозируется на эпистемическую неопределённость,
\begin{equation}
U_{\mathrm{epi}}
= \mathrm{Var}_{p(\theta|\calD)}\!\bigl[\E_{p(y|\bx,\theta)}[y]\bigr]
\approx \frac{1}{M}\sum_{m=1}^{M}(\hat{y}^{(m)}-\bar{y})^2,
\label{eq:stoch_epi}
\end{equation}
измеряющую дисперсию ожидаемого предсказания по выборкам параметров, и алеаторную неопределённость,
\begin{equation}
U_{\mathrm{ale}}
= \E_{p(\theta|\calD)}\!\bigl[\mathrm{Var}_{p(y|\bx,\theta)}[y]\bigr]
\approx \frac{1}{M}\sum_{m=1}^{M}\hat{\sigma}^{(m)2},
\label{eq:stoch_ale}
\end{equation}
усредняющую предсказанную дисперсию по выборкам параметров, где $\hat{y}^{(m)}$ --- среднее предсказание и $\hat{\sigma}^{(m)2}$ --- предсказанная дисперсия из выборки~$m$, а $\bar{y}=\frac{1}{M}\sum_m\hat{y}^{(m)}$.

\subsection{Стохастическое состязательное обучение}

Вместо генерации наихудших возмущений при фиксированных бюджетах стохастическое состязательное обучение выбирает возмущения из обучаемого распределения~\cite{Madry2018,Zhang2019trades,Shafahi2019},
\begin{equation}
\bdelta\sim q_\psi(\bdelta|\bx)
= \calN\!\bigl(\mu_\psi(\bx),\;\sigma_\psi(\bx)^2\mathbf{I}\bigr),
\label{eq:stoch_pert_dist}
\end{equation}
параметризованного нейронными сетями $\mu_\psi,\sigma_\psi$, обучающими пространственно-адаптивную статистику возмущений. Целевая функция обучения --- стохастическая задача мин-макс
\begin{equation}
\min_\theta\;\max_\psi\;
\E_{(\bx,y)\sim\calD}\;
\E_{\bdelta\sim q_\psi(\bdelta|\bx)}\!
\bigl[\calL(f_\theta(\bx+\bdelta),y)\bigr],
\label{eq:stoch_adv_obj}
\end{equation}
решаемая поочерёдным градиентным подъёмом по~$\psi$ и спуском по~$\theta$.

\subsection{Многоцелевые потери}

Полные потери обучения объединяют пять членов,
\begin{equation}
\calL_{\mathrm{total}}
= \lambda_1\calL_{\mathrm{task}}
+ \lambda_2\calL_{\mathrm{KL}}
+ \lambda_3\calL_{\mathrm{adv}}
+ \lambda_4\calL_{\mathrm{calib}}
+ \lambda_5\calL_{\mathrm{reg}},
\label{eq:stoch_total_loss}
\end{equation}
где $\calL_{\mathrm{task}}=\E_{(\bx,y)}[-\log p_\theta(y|\bx)]$ --- потери классификации, $\calL_{\mathrm{KL}}~\cite{Blundell2015}=\mathrm{KL}(q_\phi(\theta)\|p(\theta))$ регуляризует вариационное апостериорное, $\calL_{\mathrm{adv}}=\E_{(\bx,y)}\E_{\bdelta\sim q_\psi}[-\log p_\theta(y|\bx{+}\bdelta)]$ обеспечивает состязательную робастность, $\calL_{\mathrm{calib}}=\mathrm{ECE}~\cite{Guo2017,Naeini2015}(\{(p_\theta(\hat{y}|\bx_i),\mathbf{1}[\hat{y}=y_i])\}_i)$ штрафует ошибку калибровки, а $\calL_{\mathrm{reg}}=\|\theta\|_2^2$ контролирует величину весов.

\begin{algorithm}[htbp]
\caption{Стохастический трансформер: вариационное байесовское обучение с состязательной робастностью}\label{alg:stoch_trans}
\KwInput{Обучающие данные $\calD$, априорное $p(\theta)$, выборки МК $M$, параметры распределения возмущений $\psi_0$, веса потерь $\lambda_1,\ldots,\lambda_5$}
\KwOutput{Вариационные параметры $\phi^*$, параметры возмущений $\psi^*$}
Инициализировать вариационные параметры $\phi = \{\bm{\mu}_Q,\bm{\sigma}_Q,\bm{\mu}_K,\bm{\sigma}_K,\bm{\mu}_V,\bm{\sigma}_V\}$\;
Инициализировать параметры сети возмущений $\psi$\;
\For{итерация $t = 1,\ldots,T$}{
  \For{каждый мини-батч $\{(\bx_i,y_i)\}$}{
    \textbf{Шаг возмущения:}
    Вычислить $\mu_\psi(\bx_i),\sigma_\psi(\bx_i)$ и выбрать $\bdelta_i\sim\calN(\mu_\psi(\bx_i),\sigma_\psi(\bx_i)^2\mathbf{I})$\;
    Обновить $\psi \leftarrow \psi + \alpha_\psi\nabla_\psi\frac{1}{M}\sum_{m=1}^M\calL(f_{\theta^{(m)}}(\bx_i+\bdelta_i),y_i)$\;
    \textbf{Шаг модели:}
    \For{$m = 1,\ldots,M$}{
      Выбрать $\bW_Q^{(m)}=\bm{\mu}_Q+\bm{\sigma}_Q\odot\bm{\epsilon}^{(m)}$, аналогично для $\bW_K^{(m)},\bW_V^{(m)}$\;
      Вычислить выход внимания и предсказания $\hat{y}^{(m)},\hat{\sigma}^{(m)2}$\;
    }
    Вычислить $U_{\mathrm{epi}} = \frac{1}{M}\sum_m(\hat{y}^{(m)}-\bar{y})^2$, $U_{\mathrm{ale}} = \frac{1}{M}\sum_m\hat{\sigma}^{(m)2}$\;
    Вычислить $\calL_{\mathrm{total}} = \lambda_1\calL_{\mathrm{task}} + \lambda_2\calL_{\mathrm{KL}} + \lambda_3\calL_{\mathrm{adv}} + \lambda_4\calL_{\mathrm{calib}} + \lambda_5\calL_{\mathrm{reg}}$\;
    Обновить $\phi \leftarrow \phi - \alpha_\phi\nabla_\phi\calL_{\mathrm{total}}$\;
  }
}
\Return $\phi^*,\psi^*$ вместе с $U_{\mathrm{epi}},U_{\mathrm{ale}}$ для каждого тестового входа
\end{algorithm}


%%=====================================================================
\section{Прямо совместимое обнаружение при смене протокольного режима}
\label{sec:pq_idps}
%%=====================================================================

Данный раздел адресует Подзадачу~5 (Пробел~5: Нестационарность $\times$ Робастность) путём построения архитектуры обнаружения, сохраняющей робастность при структурном изменении порождающих данные протоколов графа, обеспечивая прямую совместимость, необходимую для предметно-ориентированной фундаментальной модели (CyberSecLLM).

\subsection{Протокол-специфичное извлечение признаков}

При развёртывании нового протокольного режима наблюдаемые признаки на рёбрах графа, такие как распределения времени, размеры сообщений и подсчёты взаимодействий, изменяются характерным образом. Для протоколов на основе решёток, например, тайминг операций обмена ключами демонстрирует вариации, коррелирующие с вычислительной структурой решёточных операций. Эти вариации захватываются как вектор статистических признаков,
\begin{equation}
\mathbf{f}_{\mathrm{lattice}}
= \bigl[\mu_t,\;\sigma_t,\;\mathrm{skew}(t),\;\mathrm{kurt}(t),\;
  \mathrm{pctl}(t),\;\mathrm{acf}(t)\bigr],
\label{eq:pq_feat_lattice}
\end{equation}
где $t=(t_1,\dots,t_n)$ обозначает измерения операционного тайминга, а записи захватывают центральную тенденцию, дисперсию, форму распределения и темпоральную зависимость. Аналогичные вектора признаков $\mathbf{f}_{\mathrm{code}}$ и $\mathbf{f}_{\mathrm{hash}}$ определены для семейств протоколов на основе кодов и хэшей, включая структурные параметры такие как размеры ключей, веса ошибок, энтропия синдрома, размеры подписей, подсчёты вызовов хэш-функций и глубины деревьев Меркла.

\subsection{Унифицированная мультипротокольная архитектура обнаружения}

Протокол-специфичные кодировщики формируют латентные представления $\mathbf{z}_{\mathrm{lattice}}=\mathrm{ENC}_{\mathrm{lattice}}(\mathbf{f}_{\mathrm{lattice}})$, $\mathbf{z}_{\mathrm{code}}=\mathrm{ENC}_{\mathrm{code}}(\mathbf{f}_{\mathrm{code}})$, $\mathbf{z}_{\mathrm{hash}}=\mathrm{ENC}_{\mathrm{hash}}(\mathbf{f}_{\mathrm{hash}})$. Слияние кросс-вниманием обеспечивает взаимодействие между тремя протокольными представлениями,
\begin{equation}
\mathbf{z}_{\mathrm{fused}}
= \mathrm{MultiHeadAttention}\!\bigl(
  [\mathbf{z}_{\mathrm{lattice}},\;
   \mathbf{z}_{\mathrm{code}},\;
   \mathbf{z}_{\mathrm{hash}}]\bigr),
\label{eq:pq_fusion}
\end{equation}
и многозадачная классификация предсказывает как метку категории, так и целевой протокол,
\begin{equation}
p(\mathrm{class},\mathrm{protocol})
= \mathrm{Softmax}\!\bigl(\bW_{\mathrm{out}}\mathbf{z}_{\mathrm{fused}}
  + \mathbf{b}_{\mathrm{out}}\bigr).
\label{eq:pq_class}
\end{equation}

\subsection{Формальная сводка безопасности}

Гарантия робастности системы обнаружения при новом протокольном режиме наследуется от сложности задачи обучения с ошибками (LWE)~\cite{Regev2009,Peikert2016}. При $\mathbf{A}\in\Z_q^{m\times n}$ и $\mathbf{b}=\mathbf{A}\mathbf{s}+\mathbf{e}\bmod q$ для секрета~$\mathbf{s}$ и малой ошибки~$\mathbf{e}$ любой полиномиально-временной противник, различающий легитимную и манипулированную последовательность согласования протокола, может быть использован для построения решателя LWE, что противоречит предполагаемой сложности. Формальная сводка, представленная полностью в Приложении~А, устанавливает, что ошибка второго рода классификатора ограничена преимуществом лучшего решателя LWE плюс статистический член, убывающий с числом обучающих образцов.

\begin{algorithm}[htbp]
\caption{Прямо совместимое мультипротокольное обнаружение}\label{alg:pqidps}
\KwInput{Поток графовых событий со смешанными протокольными режимами на рёбрах, протокол-специфичные кодировщики, бюджет состязат. обучения $\epsilon$}
\KwOutput{Классификация $(\mathrm{class},\mathrm{protocol})$ и сертификат робастности на основе LWE}
\For{каждое входящее событие графа}{
  Определить протокольный режим через анализ последовательности согласования\;
  \uIf{на основе решёток (Kyber/Dilithium)}{
    Извлечь признаки тайминга $\mathbf{f}_{\mathrm{lattice}} = [\mu_t,\sigma_t,\mathrm{skew},\mathrm{kurt},\mathrm{pctl},\mathrm{acf}]$\;
    Кодировать: $\mathbf{z}_{\mathrm{lattice}} \leftarrow \mathrm{ENC}_{\mathrm{lattice}}(\mathbf{f}_{\mathrm{lattice}})$\;
  }\uElseIf{на основе кодов (McEliece/Niederreiter)}{
    Извлечь структурные признаки $\mathbf{f}_{\mathrm{code}}$; кодировать $\mathbf{z}_{\mathrm{code}} \leftarrow \mathrm{ENC}_{\mathrm{code}}(\mathbf{f}_{\mathrm{code}})$\;
  }\Else{
    Извлечь хэш-признаки $\mathbf{f}_{\mathrm{hash}}$; кодировать $\mathbf{z}_{\mathrm{hash}} \leftarrow \mathrm{ENC}_{\mathrm{hash}}(\mathbf{f}_{\mathrm{hash}})$\;
  }
  Слить: $\mathbf{z}_{\mathrm{fused}} \leftarrow \mathrm{MultiHeadAttention}([\mathbf{z}_{\mathrm{lattice}},\mathbf{z}_{\mathrm{code}},\mathbf{z}_{\mathrm{hash}}])$\;
  Классифицировать: $p(\mathrm{class},\mathrm{protocol}) \leftarrow \mathrm{Softmax}(\bW_{\mathrm{out}}\mathbf{z}_{\mathrm{fused}} + \mathbf{b}_{\mathrm{out}})$\;
}
\textbf{Обучение:} Чередование стандартных прямых проходов и состязательных прямых проходов с PGD-возмущениями $\|\bdelta\|\le\epsilon$, применяемыми к векторам признаков\;
\textbf{Сертификация:} Проверить сводку LWE (Приложение~А) и сообщить параметр робастности $\lambda_{\mathrm{sec}}$\;
\Return классификацию и сертификат робастности
\end{algorithm}


%%=====================================================================
\section{Метод М7: Теоретико-игровая сертификация и единая система}
\label{sec:game}
%%=====================================================================

Данный раздел решает Подзадачу~7 путём разработки формулировки игры Штакельберга и члена связки согласованности $\sum_{i<j}\calL_{ij}(\theta_i,\theta_j)$ из единой задачи (Уравнение~\eqref{eq:ps_unified}). Вместе с компонентами сертификации это составляет Метод М7 и адресует Пробел~7 (все условия $\times$ все требования).

\subsection{Формулировка игры Штакельберга}

Защитник выбирает модель $f_\theta$ и конфигурацию развёртывания~$c$, включающую параметры политики обнаружения. Противник наблюдает $(f_\theta,c)$ и выбирает стратегию возмущения~$a$. Функции полезности защитника и противника соответственно:
\begin{equation}
U_d(f_\theta,c,a)
= \E_{\bx\sim\calD(a)}\!\bigl[\mathbf{1}[f_\theta(\bx)=y]\bigr] - \mathrm{Cost}(c),
\label{eq:game_ud}
\end{equation}
\begin{equation}
U_a(f_\theta,c,a)
= \Prob[\text{ошибка классиф.}\mid f_\theta,c,a] - \mathrm{Cost}(a),
\label{eq:game_ua}
\end{equation}
где $\calD(a)$ --- распределение данных, индуцированное стратегией противника. Равновесие Штакельберга --- решение
\begin{equation}
(f_\theta^*,c^*)
= \arg\max_{f_\theta,c}\;U_d\!\bigl(f_\theta,c,\;a^*(f_\theta,c)\bigr),
\qquad
a^*(f_\theta,c)=\arg\max_a\;U_a(f_\theta,c,a).
\label{eq:game_stackel}
\end{equation}

\subsection{Характеристика равновесия Нэша}

При одновременной игре профиль смешанных стратегий $(\pi_d,\pi_a)$ является равновесием Нэша, когда ни один игрок не улучшает полезность односторонним отклонением,
\begin{equation}
\E_{\pi_a}\!\bigl[U_d(\pi_d,a)\bigr]
\ge \E_{\pi_a}\!\bigl[U_d(\pi_d',a)\bigr]\;\forall\,\pi_d',
\qquad
\E_{\pi_d}\!\bigl[U_a(f_\theta,\pi_a)\bigr]
\ge \E_{\pi_d}\!\bigl[U_a(f_\theta,\pi_a')\bigr]\;\forall\,\pi_a'.
\label{eq:game_nash}
\end{equation}
В частном случае игры с нулевой суммой $U_d+U_a=0$ равновесие совпадает с минимаксным значением $\pi_d^*=\arg\max_{\pi_d}\min_a\E_{\pi_d}[U_d(f_\theta,a)]$.

\subsection{Онлайн-обучение с гарантиями отсутствия регрета}

На раунде~$t$ защитник выбирает конфигурацию $c_t\in\calC$ и наблюдает вознаграждение $r_t=U_d(c_t,a_t)$. Кумулятивный регрет после $T$ раундов:
\begin{equation}
R(T)
= \max_{c\in\calC}\sum_{t=1}^{T}U_d(c,a_t)
  - \sum_{t=1}^{T}U_d(c_t,a_t).
\label{eq:game_regret}
\end{equation}
Обновление мультипликативными весами~\cite{AroraHK2012,CesaBianchi2006} поддерживает распределение $p_t(c)\propto\exp\!\bigl(\eta\sum_{s=1}^{t-1}\mathbf{1}[c_s{=}c]\,r_s\bigr)$ со скоростью обучения $\eta=\sqrt{\log|\calC|/T}$.

\begin{theorem}[Граница регрета]
\label{thm:regret}
Алгоритм мультипликативных весов со скоростью обучения $\eta=\sqrt{\log|\calC|/T}$ достигает
\begin{equation}
R(T)\le 2\sqrt{T\log|\calC|},
\label{eq:game_regret_bound}
\end{equation}
так что средний регрет на раунд $R(T)/T=\calO(\sqrt{\log|\calC|/T})\to 0$ при $T\to\infty$.
\end{theorem}

\noindent\textit{Эскиз доказательства.} Доказательство проходит через ограничение потенциальной функции $\Phi_t=\log\sum_c\exp(\eta\sum_{s=1}^{t}r_s(c))$ сверху и снизу. Верхняя граница использует вогнутость логарифма и определение алгоритма; нижняя граница выделяет лучшее действие ретроспективно. Объединение двух границ и оптимизация по~$\eta$ даёт указанную скорость. Полный аргумент приведён в Приложении~А.\qed

\subsection{Выбор стратегии с учётом неопределённости}

Байесовская неопределённость стохастического трансформера из Раздела~\ref{sec:stochastic} информирует полезность защитника через выигрыш со штрафом за неопределённость,
\begin{equation}
\tilde{U}_d(c,a)
= \begin{cases}
U_d(c,a) - \kappa\,U_{\mathrm{epi}} & \text{если}\;U_{\mathrm{epi}}<\tau,\\
-\infty & \text{иначе},
\end{cases}
\label{eq:game_unc}
\end{equation}
где $\kappa$ дисконтирует полезность при неопределённости, а порог~$\tau$ инициирует направление на экспертный анализ. Это порождает чувствительные к риску политики, воздерживающиеся от предсказаний при недостаточной уверенности модели.

\begin{algorithm}[htbp]
\caption{Теоретико-игровая робастность с онлайн-обучением, учитывающим неопределённость}\label{alg:game}
\KwInput{Пространство действий $\calC$, порог неопределённости $\tau$, штраф $\kappa$, расписание скорости обучения}
\KwOutput{Политика защитника $\pi_d^*$ с гарантией регрета}
Инициализировать равномерное распределение $p_1(c) = 1/|\calC|$ для всех $c\in\calC$\;
Установить $\eta = \sqrt{\log|\calC|/T}$\;
\For{раунд $t = 1,\ldots,T$}{
  Выбрать действие защитника $c_t \sim p_t$\;
  Наблюдать действие противника $a_t$ и вознаграждение $r_t = U_d(c_t,a_t)$\;
  Запросить стохастический трансформер для эпистемической неопределённости $U_{\mathrm{epi}}$\;
  \uIf{$U_{\mathrm{epi}} \ge \tau$}{
    Установить $\tilde{r}_t = -\infty$ (отклонить; направить на экспертный анализ)\;
  }\Else{
    Установить $\tilde{r}_t = r_t - \kappa\,U_{\mathrm{epi}}$\;
  }
  Обновить веса: $p_{t+1}(c) \propto p_t(c)\exp(\eta\,\tilde{r}_t\,\mathbf{1}[c_t=c])$ для всех $c\in\calC$\;
  Нормализовать $p_{t+1}$\;
}
Вычислить эмпирическое равновесие Нэша из частот игры $\{(c_t,a_t)\}_{t=1}^T$\;
\Return $\pi_d^*$ и подтверждённый регрет $R(T)\le 2\sqrt{T\log|\calC|}$
\end{algorithm}


%%=====================================================================
\subsection{Единая сертификация через связанную оптимизацию}
\label{sec:unified}
%%=====================================================================

Семь разработанных выше методов (М1--М6, а также теоретико-игровой компонент М7) разделяют математические структуры --- анализ сопряжённых чувствительностей, геометрию Вассерштейна, регуляризацию KL-дивергенцией и динамику с ограниченной константой Липшица --- допускающие единую трактовку через связанную оптимизацию. Данный раздел завершает Метод М7 разработкой системы сертификации, интегрирующей все компонентные гарантии.

\subsection{Связанная целевая функция}

Пусть $\theta=\{\theta_{\mathrm{CT}},\theta_{\mathrm{TE}},\theta_{\mathrm{Fed}},\theta_{\mathrm{PQ}},\theta_{\mathrm{MS}},\theta_{\mathrm{ST}},\theta_{\mathrm{GT}}\}$ собирает параметры всех семи методов. Единая целевая функция связывает компонентные потери с попарными членами согласованности,
\begin{equation}
\calL_{\mathrm{unified}}(\theta)
= \sum_{i=1}^{7}\lambda_i\,\calL_i(\theta_i)
+ \sum_{i<j}\lambda_{ij}\,\calL_{ij}(\theta_i,\theta_j),
\label{eq:unified_obj}
\end{equation}
где каждые $\calL_i$ --- потери, специфичные для метода, из соответствующего раздела, а потери связки
\begin{equation}
\calL_{ij}(\theta_i,\theta_j)
= \mathrm{KL}\!\bigl(p_{\theta_i}(y|\bx)\|p_{\theta_j}(y|\bx)\bigr)
+ \|\mathbf{z}_{\theta_i}(\bx)-\mathbf{z}_{\theta_j}(\bx)\|_2^2
\label{eq:unified_coupling}
\end{equation}
обеспечивают согласованность между предсказательными распределениями и обученными представлениями методов~$i$ и~$j$.

\subsection{Анализ сходимости}

\begin{theorem}[Сходимость единой системы]
\label{thm:unified}
Предположим, что каждые компонентные потери $\calL_i$ являются $L_i$-гладкими и каждые потери связки $\calL_{ij}$ являются $L_{ij}$-гладкими. При поочерёдном градиентном спуске со скоростями обучения $\eta_i\le 1/(L_i+\sum_j L_{ij})$ единая оптимизация сходится со скоростью
\begin{equation}
\|\nabla\calL_{\mathrm{unified}}(\theta_T)\|^2
\le \frac{2\bigl(\calL_{\mathrm{unified}}(\theta_0)
  -\calL_{\mathrm{unified}}(\theta^*)\bigr)}{\eta_{\min}\,T},
\label{eq:unified_conv}
\end{equation}
где $\eta_{\min}=\min_i\eta_i$, а $\theta^*$ --- стационарная точка.
\end{theorem}

\noindent\textit{Эскиз доказательства.} Каждый шаг поочерёдного градиента по компоненту~$i$ уменьшает единую целевую функцию не менее чем на $\eta_i\|\nabla_{\theta_i}\calL_{\mathrm{unified}}\|^2/2$ при условии гладкости, поскольку члены связки гладки по каждому блоку переменных. Суммирование по всем компонентам и всем $T$ раундам, деление на $T$ и взятие минимума по скоростям обучения даёт указанную границу. Полный вывод, включая обработку перекрёстных градиентов, приведён в Приложении~А.\qed

\begin{algorithm}[htbp]
\caption{Единая система: поочерёдная оптимизация со связкой согласованности}\label{alg:unified}
\KwInput{Компонентные потери $\calL_1,\ldots,\calL_7$, веса связки $\{\lambda_{ij}\}$, константы гладкости $\{L_i,L_{ij}\}$}
\KwOutput{Совместно оптимизированные параметры $\theta^*=\{\theta_1^*,\ldots,\theta_7^*\}$}
Инициализировать все компонентные параметры $\theta_1^{(0)},\ldots,\theta_7^{(0)}$\;
Установить скорости обучения $\eta_i = 1/(L_i + \sum_j L_{ij})$ для каждого компонента\;
\For{раунд $t = 0,1,\ldots,T{-}1$}{
  \For{компонент $i = 1,\ldots,7$}{
    Вычислить компонентный градиент $\mathbf{g}_i = \nabla_{\theta_i}\calL_i(\theta_i^{(t)})$\;
    Вычислить градиенты связки $\mathbf{g}_{ij} = \nabla_{\theta_i}\calL_{ij}(\theta_i^{(t)},\theta_j^{(t)})$ для всех $j\neq i$\;
    Совокупный градиент: $\mathbf{g}_i^{\mathrm{total}} = \lambda_i\mathbf{g}_i + \sum_{j\neq i}\lambda_{ij}\mathbf{g}_{ij}$\;
    Обновить: $\theta_i^{(t+1)} \leftarrow \theta_i^{(t)} - \eta_i\,\mathbf{g}_i^{\mathrm{total}}$\;
  }
  Вычислить $\calL_{\mathrm{unified}}(\theta^{(t+1)})$ и проверить критерий сходимости\;
}
\Return $\theta^* = \theta^{(T)}$
\end{algorithm}


%%=====================================================================
\section{Итоги главы}
\label{sec:ch2_summary}
%%=====================================================================

В настоящей главе разработаны полные математические основы для семи новых методов, выведенных из единой постановки задачи (Уравнение~\eqref{eq:ps_unified}) через систематическую декомпозицию на семь подзадач. Единая задача была последовательно построена путём включения состязательной робастности (Этап~2), динамики непрерывного времени с сертификацией Липшица (Этап~3), распределённого обучения с дифференциальной приватностью (Этап~4), многоуровневого представления с защитой от забывания (Этап~5), вычислительной эффективности через модели пространства состояний (Этап~6) и нестационарной адаптации с оценкой неопределённости (Этап~7).

Метод М1 (CT-TGNN) обеспечивает динамику графов непрерывного времени с сертифицированной робастностью по Липшицу, решая Подзадачу~1 (Пробел~1). Метод М2 (FedLLM-API) достигает византийско-устойчивой федеративной оптимизации с дифференциальной приватностью и выравниванием оптимального транспорта, решая Подзадачу~2 (Пробел~2). Метод М3 (TripleE-TGNN) захватывает многоуровневую структуру через тройные вложения с упругой консолидацией весов, решая Подзадачу~3 (Пробел~3). Метод М4 (MambaShield) обеспечивает эффективный последовательный вывод с PAC-байесовскими сертификатами обобщения, решая Подзадачу~4 (Пробел~4). Архитектура прямо совместимого обнаружения адресует смену протокольных режимов, решая Подзадачу~5 (Пробел~5). Методы М5/М6 (Стохастический трансформер и UC-HGP) обеспечивают калиброванную байесовскую неопределённость при состязательном возмущении, решая Подзадачу~6 (Пробел~6). Метод М7 (Сертификация), включающий теоретико-игровую систему и связку единой оптимизации, объединяет все компонентные гарантии через регуляризацию согласованности с доказанной скоростью сходимости (Теорема~\ref{thm:unified}), решая Подзадачу~7 (Пробел~7).

Каждый метод сформулирован на уровне математических операций над графовыми структурами без привязки к какой-либо конкретной прикладной области. Теоретические гарантии --- скорости сходимости, сертификаты робастности, бюджеты приватности, границы регрета и PAC-байесовские границы обобщения --- выполняются везде, где эти методы развёрнуты. В следующей главе все методы реализуются в области графов сетевого трафика, описывая процесс адаптации и согласования, через который предметно-независимая математическая основа отображается на сущности и ограничения кибербезопасности и гибридных систем обнаружения и предотвращения сетевых вторжений.


